\documentclass[a4paper,12pt]{article}
\usepackage{ava_style}
\renewcommand{\AVAManualTitle}{Spectroscopy Manual}

\graphicspath{{figures/}{figures/spectroscopy/}}

\begin{document}

\avatitlepage
  {Spectroscopy Manual}
  {0.1}
  {16 November 2026}
  {Juan V.\ Hern\'{a}ndez Santisteban}
  {jvhs1@st-andrews.ac.uk}
  {School of Physics \& Astronomy, University of St Andrews,\\
   North Haugh, St Andrews, KY16~9SS, UK}

\tableofcontents
\clearpage

%%============================================================
\section{Introduction}
\label{sec:intro}
%%============================================================

\begin{warnbox}
\textbf{Draft --- v0.1.}  This manual describes the spectroscopic pipeline
as designed and partially implemented.  The banzai-floyds reduction stage
is under active testing.  Figures in this manual use the current Fairall\,9
and Mrk\,817 test archives.
\end{warnbox}

This manual documents the spectroscopic arm of the AGN Variability Archive (AVA):
acquisition of FLOYDS spectra from the LCO archive, reduction with the improved
banzai-floyds pipeline, spectral-line fitting, and production of emission-line
light curves for reverberation mapping.

The spectroscopy pipeline works in concert with the photometry pipeline
(see the \emph{Photometry Manual}).  Its primary science outputs are:
\begin{itemize}
  \item Flux-calibrated 1-d and 2-d spectra, stored per target per night under
        \path{<archive_root>/spectroscopy/<target>/reduced/<YYYYMMDD>/}.
  \item H$\beta$, [O\,\textsc{iii}] and continuum flux measurements forming a
        spectral time series.
  \item Combined spectroscopic + photometric light curves for RM cross-correlation.
\end{itemize}

All data paths are configured in \path{config/ava_config.yaml} under
\texttt{paths.*} and \texttt{spectroscopy.*}.  No hardcoded absolute paths exist
in the production pipeline.

\subsection{FLOYDS instrument}

FLOYDS is a cross-dispersed, low-resolution spectrograph ($R \approx 300$--600)
mounted on the 2\,m Faulkes Telescopes at Haleakal\={a} (FTN, site \texttt{ogg})
and Siding Spring (FTS, site \texttt{coj}).  It simultaneously covers
$\sim$3200--10\,000\,\AA\ in two overlapping orders:
\begin{itemize}
  \item \textbf{Blue arm:} $\sim$3200--5700\,\AA\ (order 2).
  \item \textbf{Red arm:}  $\sim$5400--10\,000\,\AA\ (order 1).
\end{itemize}
The two orders overlap at $\sim$5400--5700\,\AA; the pipeline stitches them
into a single continuous spectrum for display but keeps both for analysis.
Standard stars are observed each night for flux calibration, with arcs and lamp
flats taken in the same observing block.

%%============================================================
\section{Pipeline Overview}
\label{sec:overview}
%%============================================================

The spectroscopy pipeline runs in four stages, each implemented as a Prefect
task and callable independently:

\begin{enumerate}
  \item \textbf{Download} (\path{ava/spectroscopy/download.py}) ---
        queries the LCO archive for new FLOYDS frames (science, standards,
        arcs, flats); populates the lcodb spectra database.
  \item \textbf{Reduction} (\path{ava/spectroscopy/reduce.py}) ---
        runs the improved banzai-floyds pipeline in a dedicated Python~3.12
        environment; files QC results to \texttt{spec\_qc}.
  \item \textbf{Spectral products} (\path{ava/spectroscopy/products.py}) ---
        extracts H$\beta$, [O\,\textsc{iii}] and continuum fluxes from each
        reduced spectrum; writes the emission-line time series.
  \item \textbf{Web plots} (\path{ava/spectroscopy/webplots.py}) ---
        produces the spectral ARM plot, mean spectrum, H$\beta$ series and
        fit summary for the AVA website.
\end{enumerate}

%%============================================================
\section{Database and Archive Structure}
\label{sec:database}
%%============================================================

\subsection{Directory layout}

All data live under \texttt{paths.archive\_root} (set in
\path{config/ava_config.yaml}; defaults to \texttt{data/test\_archive}
for the sandbox).  The spectroscopy subtree is:

\begin{lstlisting}[style=shell]
<archive_root>/spectroscopy/
  raw_data/<target>/YYYYMMDD/*.fits.fz   # raw LCO frames
  <target>/
    reduced/YYYYMMDD/
      *-e91-1d.fits.fz                   # flux-calibrated 1-d spectrum
      *-e91-2d.fits.fz                   # 2-d rectified spectrum (optional)
    plots/
      lco_latest_<target>_spec.png       # most-recent spectrum
      lco_latest_<target>_spec_avg.png   # mean spectrum
      lco_latest_<target>_hbeta.png      # H-beta time series
      lco_latest_<target>_hbeta_fit.png  # H-beta fit summary
  banzai/                                # banzai-floyds work tree
    input/, processed/, db.sqlite
\end{lstlisting}

\subsection{lcodb spectra database}

Spectroscopic metadata is stored in \path{data/lcodb/spectra.sqlite}.
Key tables:

\begin{center}
\small
\begin{tabular}{lp{9cm}}
\toprule
Table & Contents \\
\midrule
\texttt{spectra} & One row per archive frame: \texttt{frame\_id},
  \texttt{data\_id}, role (science/standard), night, paths \\
\texttt{spec\_qc} & One row per reduced spectrum: QC keywords from
  banzai-floyds (\texttt{L1Q*}), product paths, pipeline version \\
\texttt{standards} & Available standard stars per night per site \\
\texttt{spec\_lines} & Per-exposure line and continuum flux fits \\
\texttt{line\_variability} & Per-target summary: mean, rms, $F_\mathrm{var}$ \\
\bottomrule
\end{tabular}
\end{center}

\subsection{Standard-star coverage}

As of November 2026 the archive contains 9\,524 science spectra across 71 AGN.
\textbf{All 9\,524 have at least one matched spectrophotometric standard observed
on the same night from the same site.}  Zero science frames lack a flux calibrator.

%%============================================================
\section{Download Stage}
\label{sec:download}
%%============================================================

\path{ava/spectroscopy/download.py} queries the LCO archive API for new FLOYDS
frames and populates the \texttt{spectra} table.  It follows the same
overlap-and-lookback logic as the photometry downloader (see the
\emph{Photometry Manual}, Section~3):

\begin{itemize}
  \item Science frames (\texttt{OBJECT = <target>}, \texttt{INSTRUME = FLOYDS}).
  \item Spectrophotometric standards (the LCO standard-star list).
  \item Associated calibration frames (arcs, lamp flats) for each observed night.
\end{itemize}

Frame matching and the download state are recorded per night in the spectra
database so that partial nights can be resumed.

%%============================================================
\section{Reduction: improved banzai-floyds}
\label{sec:reduction}
%%============================================================

\begin{notebox}
\textbf{Status: testing.}  The improved banzai-floyds pipeline is a fork of the
official LCO pipeline with improved fringe correction, sky subtraction, and flux
calibration.  It is being validated on the Fairall\,9 and Mrk\,817 archives.
\end{notebox}

\subsection{Reduction environment}

banzai-floyds requires Python~3.12 with \texttt{torch} (neural-network trace
model).  It runs as a subprocess in its own environment, whose interpreter path
is set by \texttt{spectroscopy.reduction.python} in \path{config/ava_config.yaml}.
This isolates the torch dependency from the main AVA environment.

The banzai-floyds code location is set by \texttt{paths.floyds\_pipeline} and
the sky-flat library by \texttt{spectroscopy.reduction.skyflats\_dir} in the
same config file.

\subsection{Reduction stages}

For each science target night, \path{reduce.py} orchestrates:

\begin{enumerate}
  \item \textbf{Calibration database setup} (once per site): builds the arc
        solution and fringe phase-master library from sky flats via
        \path{banzai_floyds_offline setup-db} and \path{skyflats}.
  \item \textbf{Staging}: input symlinks placed in
        \path{banzai/input/<target>_<run>/} for science + standard frames with
        their associated calibration frames.
  \item \textbf{Chain reduction}: \path{banzai_floyds_offline chain} produces:
        arc solutions, flat-fielded frames, fringe-corrected science spectra, and
        flux-calibrated 1-d/2-d products.
  \item \textbf{Filing}: science products (\texttt{-1d}, \texttt{-2d}) are moved
        to \path{<target>/reduced/<YYYYMMDD>/};
        calibrations remain in the shared \path{banzai/processed/} tree.
  \item \textbf{QC}: \texttt{L1Q*} keywords and observing conditions are logged
        to the \texttt{spec\_qc} database table.
\end{enumerate}

\begin{figure}[htbp]
  \centering
  \includegraphics[width=\textwidth]{spectroscopy/compare_archive}
  \caption{Comparison of raw and reduced spectra for Mrk\,817 (blue arm shown).
    Top: raw 2-d frame showing the cross-dispersed trace, sky lines, and
    fringe pattern.  Bottom: the banzai-floyds reduced and fringe-corrected
    1-d spectrum, flux-calibrated against the nightly standard star.  The
    improvement in fringe suppression redward of 7000\,\AA\ is the primary
    motivation for the improved pipeline.}
  \label{fig:compare-archive}
\end{figure}

\subsection{QC keywords}

The \texttt{spec\_qc} table records the full set of banzai-floyds quality
keywords per exposure:

\begin{center}
\small
\begin{tabular}{lp{8cm}}
\toprule
Keyword & Description \\
\midrule
\texttt{L1QSN1}, \texttt{L1QSN2} & S/N in blue and red arms \\
\texttt{L1QFLAG} & Overall QC flag (0 = pass) \\
\texttt{L1QSTDGR} & Standard star grade (A/B/C/F) \\
\texttt{L1QSTDDT} & Standard-to-science $\Delta t$ (hours) \\
\texttt{L1QSTDSP} & Standard-to-science sky separation \\
\texttt{AIRMASS} & Mean airmass \\
\texttt{AGFWHM} & Autoguider FWHM (arcsec) \\
\texttt{WMSCLOUD} & Cloud extinction estimate (mag) \\
\texttt{SKYMAG} & Sky brightness (mag\,arcsec$^{-2}$) \\
\bottomrule
\end{tabular}
\end{center}

%%============================================================
\section{Spectral Products}
\label{sec:products}
%%============================================================

\path{ava/spectroscopy/products.py} processes the reduced 1-d spectra in
three steps: arm extraction, emission-line fitting, and writing to the database.

\subsection{Spectral arm extraction and display}

\begin{figure}[htbp]
  \centering
  \includegraphics[width=\textwidth]{spectroscopy/lco_latest_Mrk_817_spec}
  \caption{Most-recent reduced FLOYDS spectrum of Mrk\,817 (FTS, \texttt{coj2m002}).
    The blue and red arms are shown separately; the grey shaded region marks the
    overlap zone (5400--5700\,\AA).  Key emission lines are labelled: H$\beta$
    $\lambda$4861, [O\,\textsc{iii}]\,$\lambda\lambda$4959,5007, H$\alpha$
    $\lambda$6563, [N\,\textsc{ii}]\,$\lambda$6583.  The spectrum is shown in
    the observed frame; tick marks indicate the positions of host-galaxy
    absorption features.}
  \label{fig:spec-mrk817}
\end{figure}

\begin{figure}[htbp]
  \centering
  \includegraphics[width=\textwidth]{spectroscopy/lco_latest_fairall_9_spec}
  \caption{Most-recent FLOYDS spectrum of Fairall\,9 (FTN, \texttt{ogg2m001}).
    Fairall\,9 is the reference target for AVA FLOYDS commissioning; its bright
    AGN ($V \approx 13$\,mag at peak) gives S/N $\sim$80 per resolution element
    in a 1200\,s exposure in good conditions.}
  \label{fig:spec-f9}
\end{figure}

\subsection{Mean spectrum}

The mean spectrum is the inverse-variance weighted co-addition of all good
epochs in the archive.  It serves as a template for narrow-line subtraction
and for identifying artefacts in individual exposures.

\begin{figure}[htbp]
  \centering
  \includegraphics[width=\textwidth]{spectroscopy/lco_latest_Mrk_817_spec_avg}
  \caption{Mean FLOYDS spectrum of Mrk\,817 (all epochs, all sites).  The
    shaded band is the $1\sigma$ epoch-to-epoch rms, tracing the AGN
    variability amplitude across the spectrum.  The rms is highest in the
    broad H$\beta$ and H$\alpha$ wings, confirming variable broad-line emission.}
  \label{fig:spec-avg-mrk817}
\end{figure}

\begin{figure}[htbp]
  \centering
  \includegraphics[width=\textwidth]{spectroscopy/lco_latest_fairall_9_spec_avg}
  \caption{Mean FLOYDS spectrum of Fairall\,9.  The variability rms
    (shaded) peaks strongly in the broad emission lines; the continuum rms
    at 5100\,\AA\ is the RM monitoring proxy.}
  \label{fig:spec-avg-f9}
\end{figure}

\subsection{Emission-line fitting}

\begin{figure}[htbp]
  \centering
  \includegraphics[width=0.9\textwidth]{spectroscopy/hbeta_fit_arms}
  \caption{H$\beta$-region fit for a single Fairall\,9 epoch (blue arm, 2026
    October 13).  The model (orange) comprises: narrow H$\beta$ (blue),
    broad H$\beta$ (red), [O\,\textsc{iii}]\,$\lambda\lambda$4959,5007 (green,
    constrained to a 1:3 ratio), and a linear continuum (dashed).  Residuals
    are shown in the lower panel; the fit rms is $\sim$1\,\% of the peak flux.}
  \label{fig:hbeta-fit-arms}
\end{figure}

\path{ava/spectroscopy/lines.py} fits the H$\beta$ region
($\sim$4800--4920\,\AA\ rest frame) of each spectrum:
\begin{itemize}
  \item Narrow H$\beta$ (single Gaussian, FWHM $< 500$\,km\,s$^{-1}$).
  \item Broad H$\beta$ (one or two Gaussians, FWHM $> 1000$\,km\,s$^{-1}$).
  \item [O\,\textsc{iii}]\,$\lambda\lambda$4959,5007 (doublet with fixed 1:3
        flux ratio and tied width), used as the flux-scale calibration anchor.
  \item Linear continuum under the H$\beta$+[O\,\textsc{iii}] blend.
\end{itemize}
The fit uses \texttt{scipy.optimize.curve\_fit} with weighted residuals.
The [O\,\textsc{iii}]\,5007 flux anchors the relative epoch-to-epoch flux
calibration: narrow-line variability is assumed negligible on monitoring
timescales.

\subsection{H$\beta$ time series}

\begin{figure}[htbp]
  \centering
  \includegraphics[width=\textwidth]{spectroscopy/lco_latest_Mrk_817_hbeta}
  \caption{H$\beta$ broad-line and continuum light curve for Mrk\,817
    (all epochs).  Top panel: broad H$\beta$ flux (blue), narrow H$\beta$
    flux (grey) and [O\,\textsc{iii}]\,5007 (green, expected constant).
    Middle: continuum flux at 5100\,\AA.  Bottom: S/N per epoch.
    The broad H$\beta$ clearly responds to the continuum with a lag of
    $\sim$2--6\,weeks, the primary RM signal.}
  \label{fig:hbeta-mrk817}
\end{figure}

\begin{figure}[htbp]
  \centering
  \includegraphics[width=\textwidth]{spectroscopy/lco_latest_fairall_9_hbeta}
  \caption{H$\beta$ light curve for Fairall\,9.  The [O\,\textsc{iii}] flux
    (green) is stable to $<$5\,\%, confirming the flux calibration.  The
    broad-line variability amplitude ($F_\mathrm{var} \approx 15$\,\%) is
    higher than Mrk\,817, making Fairall\,9 the preferred commissioning target
    for RM lag measurements.}
  \label{fig:hbeta-f9}
\end{figure}

\subsection{Fit summary plot}

\begin{figure}[htbp]
  \centering
  \includegraphics[width=\textwidth]{spectroscopy/lco_latest_Mrk_817_hbeta_fit}
  \caption{H$\beta$ fit summary for Mrk\,817: each panel shows one epoch's
    fit (grey: data, orange: model, blue: narrow, red: broad, green:
    [O\,\textsc{iii}]).  Epochs with S/N $<$5 or large fit residuals are
    flagged in \texttt{spec\_lines.quality}.}
  \label{fig:hbeta-fit-mrk817}
\end{figure}

\begin{figure}[htbp]
  \centering
  \includegraphics[width=\textwidth]{spectroscopy/lco_latest_fairall_9_hbeta_fit}
  \caption{H$\beta$ fit summary for Fairall\,9.  The broad H$\beta$ component
    is well-resolved at all epochs; the typical S/N in the H$\beta$ region
    is 20--60 per resolution element.}
  \label{fig:hbeta-fit-f9}
\end{figure}

\subsection{Per-exposure records}

Each fit result is written to the \texttt{spec\_lines} table:

\begin{center}
\small
\begin{tabular}{lp{8cm}}
\toprule
Column & Description \\
\midrule
\texttt{frame\_id} & Linked to \texttt{spectra.frame\_id} \\
\texttt{data\_id} & Target name \\
\texttt{MJD} & Mid-exposure Modified Julian Date \\
\texttt{f\_hbeta\_broad} & Broad H$\beta$ flux (erg\,s$^{-1}$\,cm$^{-2}$) \\
\texttt{f\_hbeta\_narrow} & Narrow H$\beta$ flux \\
\texttt{f\_oiii\_5007} & [O\,\textsc{iii}]\,5007 flux (calibration anchor) \\
\texttt{f\_cont\_5100} & Continuum flux at rest 5100\,\AA\ \\
\texttt{snr} & S/N in the H$\beta$ region \\
\texttt{quality} & Fit quality flag (0 = good) \\
\bottomrule
\end{tabular}
\end{center}

%%============================================================
\section{Reverberation Mapping Products}
\label{sec:rm}
%%============================================================

\subsection{Combined light curves}

The spectroscopic H$\beta$ flux series and the photometric $g$-band light curve
(the continuum proxy from the photometry pipeline) are aligned on a common MJD
grid.  The combined dataset is the input for RM cross-correlation analysis with
\textsc{JAVELIN} and \textsc{ICCF}.

Spectroscopic epochs are overplotted as tick marks on all photometric web plots
(see the \emph{Photometry Manual}), providing a visual check of monitoring
cadence.

\subsection{Line variability summary}

The \texttt{line\_variability} table records per-target, per-season summary
statistics for each emission line: mean flux, rms, and
$F_\mathrm{var} = \sqrt{(\sigma^2 - \delta^2)} / \langle f \rangle$ (the
excess fractional variability above measurement noise).

%%============================================================
\section{Usage}
\label{sec:usage}
%%============================================================

\begin{lstlisting}[style=shell]
# Download new FLOYDS frames for all targets
python scripts/run_pipeline.py --spec --download

# Reduce new frames with banzai-floyds for one target
python scripts/run_pipeline.py --spec --target Mrk_817

# Full spectroscopic update: download + reduce + products + plots
python scripts/run_pipeline.py --spec --all

# Read the emission-line time series
from ava.config import load_config
from ava.spectroscopy import products as P
cfg = load_config()
lc = P.hbeta_series(cfg.spec_target_dir("Mrk_817"))
\end{lstlisting}

\begin{warnbox}
The \texttt{--spec} flag activates the spectroscopic stages.  Photometric
stages run independently; the two pipelines share only the master target
list, the AVA database, and the \path{config/ava_config.yaml} configuration.
\end{warnbox}

%%============================================================
\section{Known Issues and Future Work}
\label{sec:future}
%%============================================================

\begin{enumerate}
  \item \textbf{Reduction validation ongoing.}  The improved banzai-floyds
        pipeline is being tested on the Fairall\,9 and Mrk\,817 archives.
        Production deployment is pending sign-off on the QC metrics.
  \item \textbf{Sky-subtraction model.}  The neural-network sky model requires
        the \texttt{torch} environment and site-specific sky-flat libraries.
        These are built once and cached in \path{banzai/processed/}.
  \item \textbf{Telluric correction.}  Not yet implemented; affects the red arm
        above 7000\,\AA\ (O$_2$ B band at 6900\,\AA, A band at 7600\,\AA).
  \item \textbf{H$\alpha$ fitting.}  Planned for the red arm; deferred until
        flux calibration is fully validated across both sites.
  \item \textbf{Public data release.}  Final reduced spectra and RM products
        for completed targets will be released at
        \path{https://alymantara.com/ava}.
\end{enumerate}

\end{document}
